\documentclass[11pt]{article}
\usepackage{amsmath}
\usepackage[margin=0.68in]{geometry}
\usepackage{booktabs}
\usepackage[american]{circuitikz}
\usepackage{enumitem}
\usepackage{fancyhdr}
\usepackage{graphicx}
\usepackage[table]{xcolor}
\usepackage[hidelinks]{hyperref}
\usepackage{listings}
\usepackage{microtype}
\usepackage{tabularx}
\usepackage{titlesec}
\graphicspath{{docs/lessons/assets/}{../assets/}}
\definecolor{Ink}{gray}{0}
\definecolor{CodeGray}{gray}{0.96}
\hypersetup{hidelinks,pdftitle={ADK Lesson 037: Contact Dynamics},
 pdfauthor={Kevin Shanahan},
 pdfsubject={Active-low contact qualification with a documented passive reference fixture},
 pdflang={en-US}}
\pagestyle{fancy}\fancyhf{}
\lhead{\textsc{ADK Lesson 037}}\rhead{Contact Dynamics}
\cfoot{\thepage}\setlength{\headheight}{14pt}
\setlength{\parindent}{0pt}\setlength{\parskip}{5pt}
\setlist{nosep,leftmargin=1.45em}
\titleformat{\section}{\Large\bfseries\color{Ink}}{\thesection}{0.5em}{}
\lstdefinestyle{adkcpp}{language=C++,backgroundcolor=\color{CodeGray},
 basicstyle=\ttfamily\small,keywordstyle=\color{Ink}\bfseries,
 commentstyle=\color{Ink},columns=fullflexible,keepspaces=true,
 showstringspaces=false,frame=single,rulecolor=\color{Ink},
 breaklines=true,tabsize=4}
\newcommand{\lessonbox}[1]{%
 \begin{center}\fcolorbox{Ink}{white}{\parbox{0.92\linewidth}{#1}}\end{center}}

\begin{document}
\begin{center}
 {\Huge\bfseries Contact Dynamics}\\[4pt]
 {\Large Qualify one passive contact without hiding its raw evidence}\\[8pt]
 \textit{Observe, qualify, suppress repeats, release}
\end{center}
\begin{tabularx}{\linewidth}{@{}>{\bfseries}l X >{\bfseries}l X@{}}
Lesson & 037 --- component & Time & 120--180 minutes\\
Prerequisites & Lessons 002, 006, and 007 & Board & Arduino Mega 2560\\
Evidence & TP-C and three LED witnesses & Status & host verified; bench open\\
Energy class & E1 passive contact fixture; E0 host trace separate\\
\end{tabularx}

\lessonbox{\textbf{Reference gate.} Use only a visibly identified
C\&K/Littelfuse \texttt{RB-220-07A R}, or a separately qualified
substitution. Confirm the delivered part, terminal continuity, installed
closure orientation, 10 k\(\Omega\) resistor, and literal wiring before
power. Remove USB power before wiring, moving the contact, or placing a probe.
Stop on uncertain identity, damage, unexpected continuity, the wrong
resistor, unresolved ground, or an out-of-bound prediction. The circuit and
lesson are not physical acceptance.}

\section{Mission}
\[
\text{raw level}\rightarrow\text{qualified contact}
\rightarrow\text{attack event}\rightarrow\text{release event}.
\]
\texttt{ContactDynamics} consumes copied samples. It does not measure force,
acceleration, impact energy, orientation, damage, or safety. The canonical
sketch keeps raw evidence visible while time policy rejects chatter and
suppresses repeated attacks.

\section{Inspect and orient}
\begin{center}
% ADK visual: pencil
\includegraphics[width=0.88\linewidth]{037-contact-stages-pencil.png}

\small\textit{Alternative text: pencil timeline carries an active-low raw
contact through continuous qualification, one accepted attack, refractory
suppression, a qualified release, and a held-active warning.}
\end{center}
The drawing is an orientation aid, not an electrical schematic. Match the
exact part-number marking and reject damage, uncertain identity, mercury,
unknown capsules, and bare piezo elements. The two switch terminals are
electrically nonpolar, but installed orientation determines when the rolling
ball closes them and must be recorded.

\section{Electrically authoritative reference circuit}
\begin{center}
% ADK visual: schematic
\begin{circuitikz}
\draw (0,3) node[left]{Mega 5 V}
  to[R,l=10 k$\Omega$] (2.5,3)
  -- (3.4,3) node[circ,label=above:{TP-C / D22}] (tpc) {}
  to[nos,l_={RB-220-07A R}] (6.4,3)
  -- (6.4,1.8) node[ground]{};
\draw (0,1.2) node[left]{D30}
  to[R,l=1 k$\Omega$] (2,1.2)
  to[leD,l=raw witness] (4,1.2)
  -- (4,0.5) node[ground]{};
\draw (5,1.2) node[left]{D31}
  to[R,l=1 k$\Omega$] (7,1.2)
  to[leD,l=qualified witness] (9,1.2)
  -- (9,0.5) node[ground]{};
\draw (10,1.2) node[left]{D32}
  to[R,l=1 k$\Omega$] (12,1.2)
  to[leD,l=attack witness] (14,1.2)
  -- (14,0.5) node[ground]{};
\end{circuitikz}

\small\textit{Formal schematic: the named manufacturer contact closes TP-C
to ground; the external pull-up, not an internal pull, establishes High when
open. Each evidence LED has its own current-limiting resistor.}
\end{center}
At nominal 5 V, closure current is \(5\text{ V}/10\text{ k}\Omega=0.5\)
mA. This lies between the datasheet's 10 \(\mu\)A minimum rated-load point
and 1 mA maximum resistive load. E1 must measure the actual rail and resistor
tolerance and prove worst-case current no greater than 1 mA.

\begin{tabularx}{\linewidth}{@{}l l X@{}}
\toprule
Evidence & Mega pin & Meaning\\
\midrule
TP-C & D22 & High when open; Low when the reference contact closes\\
raw LED & D30 through 1 k\(\Omega\) & sampled contact is closed\\
accepted LED & D31 through 1 k\(\Omega\) & 150 ms witness for an accepted attack\\
ready/fault LED & D32 through 1 k\(\Omega\) & acquired and healthy; dark on stuck/fault\\
\bottomrule
\end{tabularx}
Configure \texttt{DigitalInput} with \texttt{Pull::None}; configure
\texttt{ContactDynamics} active-low.

\section{Policy contract}
Construction is inert. \texttt{initialize()} validates nonzero,
half-range-safe durations. \texttt{update(sample)} processes one copied
level, timestamp, and source status. \texttt{snapshot()} is stable and
non-consuming; \texttt{reset()} clears history and counters.

An attack requires continuous activity through \texttt{qualify}; release
requires continuous inactivity through \texttt{release}. Refractory begins at
an accepted attack. Completion before its exact boundary is suppressed, while
completion at the boundary is accepted. A held contact creates no repeated
attacks and reaches \texttt{StuckActive} at its configured boundary.

Source and timing faults create no event and require reset. Natural 32-bit
rollover is valid. A half-range or larger apparent jump is invalid. An
identical same-time sample is idempotent; changed same-time evidence faults.

\section{Narrative Mega example}
\begin{lstlisting}[style=adkcpp]
void loop ()
{
    const adk::TimePoint now (millis ());

    if (now.elapsedSince (startedAt).milliseconds () >=
            observationDurationMs
        || !observeContact (now) || !qualifyDynamics ()
        || !presentContactEvidence (now) || !actuateEvidence ())
    {
        stopSafely ();
    }
}
\end{lstlisting}
The sketch acquires D22 and all three LEDs, configures a bounded run, shows
ready, then preserves observe--qualify--present--actuate order. Any failure and the
120-second boundary call \texttt{stopSafely()}. Serial is unnecessary.

\section{Absent-hardware trace laboratory}
\begin{center}
% ADK visual: pencil
\includegraphics[width=0.82\linewidth]{037-trace-lab-pencil.png}

\small\textit{Alternative text: four pencil laboratory cards show a short
pulse rejected before qualification, chatter settling into one attack, a
second attack suppressed during refractory, and a held contact warning.}
\end{center}
No physical contact is needed. Use active-low policy with 30 ms qualification,
20 ms release, 80 ms refractory, and 120 ms stuck-active time.

\begin{tabularx}{\linewidth}{@{}r l X@{}}
\toprule Time (ms) & Level & Prediction\\
\midrule
0 & High & idle; counts 0/0\\
10 & Low & active candidate starts\\
20 & High & chatter cancels candidate\\
30 & Low & candidate restarts\\
60 & Low & accepted attack; counts 1/0\\
70 & High & release candidate\\
90 & High & release; pulse width 30 ms\\
100 & Low & candidate during refractory\\
130 & Low & suppressed attack; counts 1/1\\
\bottomrule
\end{tabularx}
Save complete status and snapshot fields after every row. Reset, replay, and
compare fields rather than C++ object bytes. This establishes deterministic
policy only; sampling cannot reveal a pulse entirely between polls.

Run a separate shortened E0 held-contact trace with
\texttt{stuckActive=120 ms}: \texttt{(0,High)}, \texttt{(10,Low)},
\texttt{(40,Low)}, \texttt{(160,Low)}, \texttt{(170,High)}, and
\texttt{(190,High)}. Predict acceptance at 40 ms, \texttt{StuckActive}
exactly at 160 ms, and release with 150 ms pulse width at 190 ms. The physical
sketch deliberately uses a longer 2 s boundary.

\section{Predict, observe, interpret}
\begin{enumerate}
\item Before power, predict open and closed TP-C levels, all three LEDs, and
closed current. Inspect the part and resistor; unresolved identity stops.
\item Observe TP-C relative to Mega ground, then raw, accepted, and ready/fault
LEDs. Record time and do not revise the prediction.
\item Interpret a discrepancy by separating circuit level, sampling, and
policy. Do not infer physical intensity from an accepted attack.
\end{enumerate}
Acquisition and safe state are separate evidence. Successful start proves that
D22 and the three output claims were acquired. Bench work must independently
measure that shutdown releases every pin to high impedance.

\section{Diagnosis}
\begin{tabularx}{\linewidth}{@{}X X@{}}
\toprule Observation & Inspect next\\
\midrule
TP-C never changes & part identity, continuity, orientation, ground path\\
raw LED disagrees with TP-C & D22 mapping, active level, sample timing\\
accepted LED never lights & continuous duration and chatter\\
several accepted pulses from one hold & refractory and held-contact policy\\
all witnesses go dark & acquisition, source status, bounded stop\\
\bottomrule
\end{tabularx}

\section{Open bench acceptance}
Record the exact delivered part, board revision, supply and instrument,
contact-resistor measurement, each LED resistor and current bound, predicted
levels/current, observed TP-C and LED states, worst-case closed current, and
shutdown high impedance. Leave every row blank
until a human performs it. Compilation, trace replay, and this PDF are not
physical verification.

\section{References}
\begin{itemize}
\item C\&K/Littelfuse, \href{https://www.ckswitches.com/media/1317/rb.pdf}
{RB series manufacturer datasheet}.
\item \texttt{src/contact\_dynamics.h} and
\texttt{tests/test\_contact\_dynamics.cpp}.
\item \texttt{examples/Lesson037ContactDynamics/}.
\end{itemize}

\lessonbox{\textbf{Result.} Host replay and Mega compilation establish the
software path for the documented reference design. Incoming conformance,
powered characterization, current measurement, safe-state measurement, and
bench acceptance remain open.}
\end{document}
